% Options for packages loaded elsewhere
\PassOptionsToPackage{unicode}{hyperref}
\PassOptionsToPackage{hyphens}{url}
\PassOptionsToPackage{dvipsnames,svgnames,x11names}{xcolor}
%
\documentclass[
  12pt,
  a4paper,
]{article}
\usepackage{amsmath,amssymb}
\usepackage{setspace}
\usepackage{iftex}
\ifPDFTeX
  \usepackage[T1]{fontenc}
  \usepackage[utf8]{inputenc}
  \usepackage{textcomp} % provide euro and other symbols
\else % if luatex or xetex
  \usepackage{unicode-math} % this also loads fontspec
  \defaultfontfeatures{Scale=MatchLowercase}
  \defaultfontfeatures[\rmfamily]{Ligatures=TeX,Scale=1}
\fi
\usepackage{lmodern}
\ifPDFTeX\else
  % xetex/luatex font selection
    \setmainfont[]{Times New Roman}
\fi
% Use upquote if available, for straight quotes in verbatim environments
\IfFileExists{upquote.sty}{\usepackage{upquote}}{}
\IfFileExists{microtype.sty}{% use microtype if available
  \usepackage[]{microtype}
  \UseMicrotypeSet[protrusion]{basicmath} % disable protrusion for tt fonts
}{}
\makeatletter
\@ifundefined{KOMAClassName}{% if non-KOMA class
  \IfFileExists{parskip.sty}{%
    \usepackage{parskip}
  }{% else
    \setlength{\parindent}{0pt}
    \setlength{\parskip}{6pt plus 2pt minus 1pt}}
}{% if KOMA class
  \KOMAoptions{parskip=half}}
\makeatother
\usepackage{xcolor}
\usepackage[margin=0.75in]{geometry}
\usepackage{color}
\usepackage{fancyvrb}
\newcommand{\VerbBar}{|}
\newcommand{\VERB}{\Verb[commandchars=\\\{\}]}
\DefineVerbatimEnvironment{Highlighting}{Verbatim}{commandchars=\\\{\}}
% Add ',fontsize=\small' for more characters per line
\usepackage{framed}
\definecolor{shadecolor}{RGB}{248,248,248}
\newenvironment{Shaded}{\begin{snugshade}}{\end{snugshade}}
\newcommand{\AlertTok}[1]{\textcolor[rgb]{0.94,0.16,0.16}{#1}}
\newcommand{\AnnotationTok}[1]{\textcolor[rgb]{0.56,0.35,0.01}{\textbf{\textit{#1}}}}
\newcommand{\AttributeTok}[1]{\textcolor[rgb]{0.13,0.29,0.53}{#1}}
\newcommand{\BaseNTok}[1]{\textcolor[rgb]{0.00,0.00,0.81}{#1}}
\newcommand{\BuiltInTok}[1]{#1}
\newcommand{\CharTok}[1]{\textcolor[rgb]{0.31,0.60,0.02}{#1}}
\newcommand{\CommentTok}[1]{\textcolor[rgb]{0.56,0.35,0.01}{\textit{#1}}}
\newcommand{\CommentVarTok}[1]{\textcolor[rgb]{0.56,0.35,0.01}{\textbf{\textit{#1}}}}
\newcommand{\ConstantTok}[1]{\textcolor[rgb]{0.56,0.35,0.01}{#1}}
\newcommand{\ControlFlowTok}[1]{\textcolor[rgb]{0.13,0.29,0.53}{\textbf{#1}}}
\newcommand{\DataTypeTok}[1]{\textcolor[rgb]{0.13,0.29,0.53}{#1}}
\newcommand{\DecValTok}[1]{\textcolor[rgb]{0.00,0.00,0.81}{#1}}
\newcommand{\DocumentationTok}[1]{\textcolor[rgb]{0.56,0.35,0.01}{\textbf{\textit{#1}}}}
\newcommand{\ErrorTok}[1]{\textcolor[rgb]{0.64,0.00,0.00}{\textbf{#1}}}
\newcommand{\ExtensionTok}[1]{#1}
\newcommand{\FloatTok}[1]{\textcolor[rgb]{0.00,0.00,0.81}{#1}}
\newcommand{\FunctionTok}[1]{\textcolor[rgb]{0.13,0.29,0.53}{\textbf{#1}}}
\newcommand{\ImportTok}[1]{#1}
\newcommand{\InformationTok}[1]{\textcolor[rgb]{0.56,0.35,0.01}{\textbf{\textit{#1}}}}
\newcommand{\KeywordTok}[1]{\textcolor[rgb]{0.13,0.29,0.53}{\textbf{#1}}}
\newcommand{\NormalTok}[1]{#1}
\newcommand{\OperatorTok}[1]{\textcolor[rgb]{0.81,0.36,0.00}{\textbf{#1}}}
\newcommand{\OtherTok}[1]{\textcolor[rgb]{0.56,0.35,0.01}{#1}}
\newcommand{\PreprocessorTok}[1]{\textcolor[rgb]{0.56,0.35,0.01}{\textit{#1}}}
\newcommand{\RegionMarkerTok}[1]{#1}
\newcommand{\SpecialCharTok}[1]{\textcolor[rgb]{0.81,0.36,0.00}{\textbf{#1}}}
\newcommand{\SpecialStringTok}[1]{\textcolor[rgb]{0.31,0.60,0.02}{#1}}
\newcommand{\StringTok}[1]{\textcolor[rgb]{0.31,0.60,0.02}{#1}}
\newcommand{\VariableTok}[1]{\textcolor[rgb]{0.00,0.00,0.00}{#1}}
\newcommand{\VerbatimStringTok}[1]{\textcolor[rgb]{0.31,0.60,0.02}{#1}}
\newcommand{\WarningTok}[1]{\textcolor[rgb]{0.56,0.35,0.01}{\textbf{\textit{#1}}}}
\usepackage{longtable,booktabs,array}
\usepackage{calc} % for calculating minipage widths
% Correct order of tables after \paragraph or \subparagraph
\usepackage{etoolbox}
\makeatletter
\patchcmd\longtable{\par}{\if@noskipsec\mbox{}\fi\par}{}{}
\makeatother
% Allow footnotes in longtable head/foot
\IfFileExists{footnotehyper.sty}{\usepackage{footnotehyper}}{\usepackage{footnote}}
\makesavenoteenv{longtable}
\usepackage{graphicx}
\makeatletter
\newsavebox\pandoc@box
\newcommand*\pandocbounded[1]{% scales image to fit in text height/width
  \sbox\pandoc@box{#1}%
  \Gscale@div\@tempa{\textheight}{\dimexpr\ht\pandoc@box+\dp\pandoc@box\relax}%
  \Gscale@div\@tempb{\linewidth}{\wd\pandoc@box}%
  \ifdim\@tempb\p@<\@tempa\p@\let\@tempa\@tempb\fi% select the smaller of both
  \ifdim\@tempa\p@<\p@\scalebox{\@tempa}{\usebox\pandoc@box}%
  \else\usebox{\pandoc@box}%
  \fi%
}
% Set default figure placement to htbp
\def\fps@figure{htbp}
\makeatother
\setlength{\emergencystretch}{3em} % prevent overfull lines
\providecommand{\tightlist}{%
  \setlength{\itemsep}{0pt}\setlength{\parskip}{0pt}}
\setcounter{secnumdepth}{5}
% definitions for citeproc citations
\NewDocumentCommand\citeproctext{}{}
\NewDocumentCommand\citeproc{mm}{%
  \begingroup\def\citeproctext{#2}\cite{#1}\endgroup}
\makeatletter
 % allow citations to break across lines
 \let\@cite@ofmt\@firstofone
 % avoid brackets around text for \cite:
 \def\@biblabel#1{}
 \def\@cite#1#2{{#1\if@tempswa , #2\fi}}
\makeatother
\newlength{\cslhangindent}
\setlength{\cslhangindent}{1.5em}
\newlength{\csllabelwidth}
\setlength{\csllabelwidth}{3em}
\newenvironment{CSLReferences}[2] % #1 hanging-indent, #2 entry-spacing
 {\begin{list}{}{%
  \setlength{\itemindent}{0pt}
  \setlength{\leftmargin}{0pt}
  \setlength{\parsep}{0pt}
  % turn on hanging indent if param 1 is 1
  \ifodd #1
   \setlength{\leftmargin}{\cslhangindent}
   \setlength{\itemindent}{-1\cslhangindent}
  \fi
  % set entry spacing
  \setlength{\itemsep}{#2\baselineskip}}}
 {\end{list}}
\usepackage{calc}
\newcommand{\CSLBlock}[1]{\hfill\break\parbox[t]{\linewidth}{\strut\ignorespaces#1\strut}}
\newcommand{\CSLLeftMargin}[1]{\parbox[t]{\csllabelwidth}{\strut#1\strut}}
\newcommand{\CSLRightInline}[1]{\parbox[t]{\linewidth - \csllabelwidth}{\strut#1\strut}}
\newcommand{\CSLIndent}[1]{\hspace{\cslhangindent}#1}
\usepackage{pdflscape}
\usepackage{caption}
\captionsetup{margin=10pt,font=small,labelfont=bf}
%\usepackage{endfloat}

\usepackage{booktabs}
\usepackage{longtable}
\usepackage{array}
\usepackage{multirow}
\usepackage{wrapfig}
\usepackage{float}
\usepackage{colortbl}
\usepackage{pdflscape}
\usepackage{tabu}
\usepackage{threeparttable}
\usepackage{threeparttablex}
\usepackage[normalem]{ulem}
\usepackage{makecell}
\usepackage{xcolor}
\usepackage{bookmark}
\IfFileExists{xurl.sty}{\usepackage{xurl}}{} % add URL line breaks if available
\urlstyle{same}
\hypersetup{
  pdftitle={Fairness or Self-Interest? An Experiment on How Communication Shapes European Perceptions of Fairness in the EU Asylum Relocation System},
  pdfauthor={(Alphabetical Order); Andrea Blasco; Begoña Cabeza Martínez; Rossella Icardi; Michal Krawczyk},
  colorlinks=true,
  linkcolor={Maroon},
  filecolor={Maroon},
  citecolor={blue},
  urlcolor={Blue},
  pdfcreator={LaTeX via pandoc}}

\title{Fairness or Self-Interest? An Experiment on How Communication Shapes European Perceptions of Fairness in the EU Asylum Relocation System}
\author{(Alphabetical Order) \and Andrea Blasco\footnote{European Commission, Joint Research Centre, Brussels, Belgium} \and Begoña Cabeza Martínez\footnote{European Commission, Joint Research Centre, Ispra, Italy} \and Rossella Icardi\footnote{European Commission, Joint Research Centre, Ispra, Italy} \and Michal Krawczyk\footnote{European Commission, Joint Research Centre, Brussels, Belgium}}
\date{This version: Jun 03, 2026}

\begin{document}
\maketitle
\begin{abstract}
Despite intense political debate over reforming the EU asylum system, there is little evidence on what drives public support for alternative solutions. Using an experimental survey, we examine Europeans' fairness perceptions of three mechanisms for distributing asylum seekers among EU countries: no relocation, relocation proportional to GDP, and relocation proportional to the resident populations. By varying the information provided about the consequences of these mechanisms, we explored the extent to which public perceptions are driven by concerns about outcomes versus principles of fairness. We find that most respondents (76\%) consider relocation to be the fairest option. However, there are large cross-country differences, with proportional relocation schemes perceived as fair by 58\% of respondents in Germany versus 87\% in Greece, which suggests a ``self-serving'' bias --- a tendency to rank as fairer the method that results in fewer asylum seekers for their own country. Moreover, when we communicate the expected number of asylum seekers for each relocation method, the self-serving bias intensifies significantly. By varying how the information is presented, we find a stronger effect when the communication displays \emph{absolute} figures rather than figures \emph{relative} to the population. These findings shed new light on how public perceptions of the fairness of the EU asylum system can be shaped by information.
\end{abstract}

{
\hypersetup{linkcolor=}
\setcounter{tocdepth}{2}
\tableofcontents
}
\setstretch{1.15}
\section{Introduction}\label{introduction}

Over the past years, the EU's asylum and migration system has been severely tested, starting with the 2015 escalation of the Syrian war and continuing through more recent crises, such as Russia's 2022 invasion of Ukraine. These events exposed a disproportionate pressure placed on Member States at the EU's external borders. In response, the EU adopted the Pact on Migration and Asylum in May 2024, which establishes a permanent solidarity framework among Member States.

Under the Pact, the country of first entry remains responsible for asylum applications, while \emph{all} Member States commit to support the system through solidarity contributions, including the relocation of applicants or beneficiaries, financial contributions, or other operational support.\footnote{See European Commission (\citeproc{ref-eu2024pact}{2024})} However, there is no consensus among countries on how this responsibility should be shared. Because the implementation of solidarity mechanisms ultimately depends on political support at the national level, understanding public perceptions of fairness and solidarity in the EU asylum and migration system is crucial for ensuring the solidarity framework functions effectively and maintaining social cohesion across the EU.

To explore these perceptions, we conducted an experimental survey in eight EU countries. Respondents were asked to evaluate the fairness of alternative asylum-seeker relocation systems under different information conditions about the consequences of each mechanism. Specifically, the survey asked respondents to rank the fairness of three mechanisms: no relocation, allocation proportional to each country's GDP, and allocation proportional to each country's population.

As in Bansak et al. (\citeproc{ref-bansak2017europeans}{2017}), we compared responses from those who received information about the consequences of implementing each method for the respondent's country against those with no such information. Extending this approach, we additionally vary how consequences are communicated, presenting them either in \emph{absolute} terms or \emph{relative} to the country's population. This distinction is important because absolute figures may exaggerate perceived burdens, while relative measures better capture per-capita responsibility and may lead to different fairness judgments.

Our work builds on a large body of literature showing that people tend to combine several fairness principles, such as need, deservingness, or reciprocity, when judging distributional outcomes (\citeproc{ref-konow2003fairest}{Konow, 2003}). However, when one's own stakes are salient, a self-serving bias can shift those judgements so that the option that benefits agents the most appears ``fairest'' (\citeproc{ref-babcock1995biased}{Babcock et al., 1995}; \citeproc{ref-konow2000fair}{Konow, 2000}). This tension between ``distributional fairness'' and ``self-serving tendency'' appears particularly relevant when considering international solidarity in the asylum system. On the one hand, the principle of fairness suggests that countries should share migrants proportionally to their capacity. On the other hand, self-serving biases may lead to prioritising national interest. However, despite pioneering work by Bansak et al. (\citeproc{ref-bansak2017europeans}{2017}) to understand this issue, how these views shape public perceptions of fairness in the migration system is still poorly understood.

Our analysis indicates that nearly three-quarters of participants perceive some form of within-EU relocation as a fair approach. However, there are significant country differences showing a tendency for participants to perceive as fairer relocation methods that result in fewer asylum seekers for their own country (``self-serving bias''). Alternative communication strategies can significantly moderate this tendency. Specifically, this study shows that presenting the likely consequences of different relocation rules in absolute terms --- reporting how many asylum seekers are expected in the participants' country per year --- significantly reinforces the self-serving bias. In contrast, presenting the same information in relative terms, rescaled by the residents' population (i.e., asylum seekers per 100,000 inhabitants), appears to reduce this bias.

\section{Related Literature}\label{related-literature}

The EU Pact on Migration and Asylum is designed to manage the complex issue of migration more effectively. Its key component is the relocation system, aiming to create a fairer and more predictable process than the previous rules. The relocation of asylum seekers has been subject to increased controversy. This issue has fuelled significant political and social tensions for several reasons, including an uneven distribution whereby border countries bear most responsibility for processing asylum claims. Other EU member states were often unwilling to accept a share of these asylum seekers, leading to a major imbalance in burden-sharing (\citeproc{ref-thielemann2018refugee}{Thielemann, 2018}). The issue has become a powerful tool for populist and far-right political parties across Europe, using anti-immigrant rhetoric to mobilise voters and gain political power (\citeproc{ref-dinas2019waking}{Dinas et al., 2019}; \citeproc{ref-halla2017immigration}{Halla et al., 2017}).

Despite an intense political debate about asylum-seekers' relocation, there is still little evidence on Europeans' preferences. Survey evidence indicates that Europeans favour capacity-based relocation over both the Dublin ``first-entry'' rule and financial transfers. The seminal paper of Bansak et al. (\citeproc{ref-bansak2017europeans}{2017}), using a survey with 18,000 respondents across 15 EU+ countries, shows that 72\% of respondents backed proportional allocation even after learning it would raise their own country's intake. Similarly, a 26-country experiment by C. W. Cappelen et al. (\citeproc{ref-cappelen2023trade}{2023}) showed that most people prefer admitting asylum seekers to paying host states. A cross-national conjoint study in eight EU countries found that citizens favour refugee protection paired with clear limits or conditions, rather than pure first-entry allocations (\citeproc{ref-jeannet2021asylum}{Jeannet et al., 2021}). Existing research, hence, suggests that public preferences are more nuanced and conditional than commonly assumed. Recent UK-based findings by Corre and Tilley (\citeproc{ref-corre2024extent}{2024}) particularly highlight this complexity. They used a conjoint experiment to test whether current UK asylum policies align with public desire. Their results challenge the common belief that the public uniformly favours stricter policies than those currently in place. These results strongly suggest that to truly understand public preferences, researchers must disaggregate complex immigration policies into their specific components, moving beyond treating asylum as a single issue.

This debate can be framed within the discussion of the Equity Theory, originally formulated by Adams (\citeproc{ref-adams1965inequity}{1965}). This framework posits that individuals evaluate the fairness of a situation by comparing the ratio of their own inputs and outcomes to those of others. In the context of asylum seekers redistribution, this translates into a demand for ``proportionality,'' where the perceived burden must be balanced against a country's capacity or resources to maintain social and political equilibrium (\citeproc{ref-thielemann2003between}{Thielemann, 2003}). When a distribution method is perceived as inequitable -e.g.~when one country feels it is bearing a disproportionate share of the responsibility relative to its neighbours- it creates a state of distress, leading to policy resistance or non-compliance (\citeproc{ref-walster1978equity}{Walster et al., 1978}). Consequently, policy outcomes in the European Union often reflect an attempt to restore this perceived equity through the implementation of ``distribution keys'' based on objective criteria such as GDP and population size (\citeproc{ref-eu2024pact}{European Commission, 2024}). These mechanisms aim to align individual country preferences with collective outcomes.

Citizens' widespread support for distributing asylum seekers based on a country's capacity implies a belief in the principle of proportional fairness and the equitable sharing of responsibility. However, this does not rule out that self-interest also shapes their attitudes, as some studies suggest. Bansak et al. (\citeproc{ref-bansak2017europeans}{2017}) found that the large majority of citizens favour an allocation policy proportional to each country's capacity over the status quo of allocation based on the country of first entry. This majority support, though weakened, persists even when respondents are informed that proportional allocation would increase the number of asylum seekers sent to their own country. These findings suggest that citizens care deeply about the fairness of the process.

Heizmann and Ziller (\citeproc{ref-heizmann2020willing}{2020}) analyses Eurobarometer data and finds that people living in countries that have taken in more asylum seekers and immigrants in the past are especially supportive of policies that redistribute wealth or resources, a preference they link to a mix of self-interest, concerns for their country's collective well-being, and moral values. Experimental work in Norway and Israel demonstrates that establishing clear proportional-allocation rules can mitigate NIMBYism (acronym for the ``not in my back yard'' attitude) by framing redistribution of refugees as a civic duty rather than an imposed burden (\citeproc{ref-solodoch2024overburdened}{Solodoch, 2024}). U.S. evidence from conjoint analysis further confirms that admission decisions are based more on collective good and ethical standards than on individual economic gain (\citeproc{ref-hainmueller2015hidden}{Hainmueller and Hopkins, 2015}).

Furthermore, a growing body of research confirms that public attitudes toward migration policy are malleable and can be meaningfully shifted through the use of concrete information and tailored framing. Studies show that specific informational framings can influence policy support: for example, informing U.S. respondents about the administrative burdens of legal immigration boosts support for more open policies (\citeproc{ref-kustov2025immigration}{Kustov and Landgrave, 2025}). Similarly, messaging congruent with personal values (using the Schwartz (\citeproc{ref-schwartz1992universals}{1992}) framework) proves more persuasive than generic humanitarian appeals in European campaigns (\citeproc{ref-dennison2020basic}{Dennison, 2020}), and correcting misperceptions about immigration's labour-market impact increases and sustains public support (\citeproc{ref-haaland2020labor}{Haaland and Roth, 2020}). In addition, narratives based on demographic challenges---such as focusing on labour shortages or an ageing population---can raise pro-immigration attitudes and political action (\citeproc{ref-facchini2022countering}{Facchini et al., 2022}). Conversely, Bansak et al. (\citeproc{ref-bansak2017europeans}{2017}) found that providing consequences-focused information (specifically, revealing that a proportional policy would increase asylum seeker numbers in a respondent's own country) appeared to boost self-interest, thereby slightly weakening support for the fairness-based proportional relocation in those countries where proportionality and self-interest don't align. A. W. Cappelen et al. (\citeproc{ref-cappelen2023second}{2023}) observed a self-serving bias where respondents deemed an allocation method fairer if it resulted in a lower national intake.

Drawing on the literature presented, we posit two core expectations (see the AEA RCT Registry for details). First, support for a proportional allocation system will be high: in each treatment, the share of participants considering the ``no relocation'' status quo as fairest will be significantly below the benchmark of one-third of all participants. Second, the perceived fairness of an option will negatively depend on the expected number of asylum seekers in one's own country under this option communicated to participants (self-serving bias). Specifically, consistent with research showing people overestimate migrant populations (\citeproc{ref-alesina2023immigration}{Alesina et al., 2023}), we expect self-serving bias to be more pronounced when the consequences are presented in absolute terms (thousands) compared to relative terms (hundreds per 100,000 inhabitants). This means that where a given relocation policy reduces the national burden with respect to the status quo (which is true for countries currently receiving a disproportionately high number of asylum applications), we predict that the policy will be rated as significantly fairer in the information treatments. Conversely, in countries in which a proportional policy increases the national burden, we predict that the policy will be rated as significantly less fair in the information treatments. In all cases, we expect that the detrimental or beneficial self-serving effect will be strongest when the information is presented in absolute terms rather than relative terms. This prediction is based on the literature suggesting that the same social effect presented in absolute rather than relative terms might be perceived as larger, affecting policy preferences (\citeproc{ref-pedersen2017ratio}{Pedersen, 2017}), a manifestation of the more general phenomenon of denominator neglect (\citeproc{ref-reyna2008numeracy}{Reyna and Brainerd, 2008}).

This work contributes to the literature on how public perceptions of fairness influence policy acceptance. Policy acceptance across all sectors depends on citizens viewing them as fair. Fairness judgments about who deserves what, often guided by deservingness heuristics, are a central component of distributive justice and shape how citizens evaluate who should contribute and who should be supported under a given policy (\citeproc{ref-almaas2010fairness}{Almås et al., 2010}; \citeproc{ref-cappelen2023trade}{C. W. Cappelen et al., 2023}; \citeproc{ref-corneo2002individual}{Corneo and Grüner, 2002}; \citeproc{ref-durante2014preferences}{Durante et al., 2014}; \citeproc{ref-fong2001social}{Fong, 2001}; \citeproc{ref-fong2011fairness}{Fong and Luttmer, 2011}; \citeproc{ref-van2006making}{Van Oorschot, 2006}, \citeproc{ref-van2000should}{2000}). Information can influence support, but its effects are heterogeneous: mobility information raises support (especially for opportunity policies), while inequality information mostly shifts concern and has large policy effects chiefly for the estate tax (\citeproc{ref-alesina2023immigration}{Alesina et al., 2023}; \citeproc{ref-kuziemko2015elastic}{Kuziemko et al., 2015}). In climate policy, support for carbon pricing rises when revenues are allocated in ways seen as fair, and acceptance can be improved with clear and trustworthy information on the policy's impact (\citeproc{ref-carattini2019win}{Carattini et al., 2019}; \citeproc{ref-maestre2019perceived}{Maestre-Andrés et al., 2019}).

Our study also adds to the literature in behavioural economics and decision making, which aims to understand how people respond to different ways of framing information. Work in medical research, for instance, has shown that people respond differently to risks when they are presented in absolute versus relative terms (\citeproc{ref-malenka1993framing}{Malenka et al., 1993}). Our work bridges these insights from medical and behavioural science literature and embeds them within a migration policy context.

\section{Methods and Data}\label{methods-and-data}

\subsection{The survey}\label{the-survey}

A sample representative of the adult population aged 18 years and over in eight EU countries (Bulgaria, France, Germany, Greece, Italy, Poland, Spain and Sweden) was drawn from Ipsos online panel by means of stratified random probability sampling. The selected countries were chosen to ensure a balanced representation of the European Union's diversity, with particular attention to coverage across macro-regions (Northern, Southern, Eastern, and Western Europe). The selection also reflects variation in key structural characteristics, including economic conditions (GDP per capita), population size, and the share of foreign-born residents.

The sampling was stratified by age, gender and geographic region (NUTS 1). Pre-defined sub-sample sizes (``quotas'') were based on official population statistics provided by Eurostat. The overall sample size of the project included 8,122 participants. The main fieldwork took place between 30 September 2024 and 21 October 2024. Details of the field work and survey design are available elsewhere (see \citeproc{ref-blasco2025fairness}{Blasco et al., 2025}).

\subsection{Experimental design}\label{experimental-design}

The survey embededed an experiment testing how different types of information affected participants' ranking of the fairness of three relocation rules for asylum seekers across EU countries:\footnote{
  The experimental design and pre-analysis plan were pre-registered within the AEA RCT Registry (AEARCTR-0014194). The Ethics Research Board of the JRC approved the experimental protocol.}

\begin{itemize}
\tightlist
\item
  \textbf{No relocation} (the baseline): asylum seekers remain in the country where they first arrive;
\item
  Allocation \textbf{proportional to each country's population}: relocations ensure equity with population as a proxy for capacity;
\item
  Allocation \textbf{proportional to each country's GDP}: relocations ensure equity with GDP as a proxy for capacity.
\end{itemize}

Respondents were asked to rank these allocation rules from the fairest (rank = 1) to the least fair (rank = 3). Before they did so, we varied the format and amount of information available across three conditions:

\begin{itemize}
\tightlist
\item
  \textbf{No information} (the control),
\item
  \textbf{Absolute} information: the expected number of asylum applications under each relocation scheme (expressed in thousands),
\item
  \textbf{Relative} information: the expected number of asylum applications under each relocation scheme per 100,000 citizens.
\end{itemize}

To illustrate the underlying relocation problem captured by our experiment, Table \ref{tab:relocation} reports the expected asylum applications in each country under each relocation rule. These figures are based on national statistics about each country's population, GDP, and historical (2014-2023) first-time asylum applications from Eurostat.\footnote{
  In this study, GDP and population are intended as proxies for a country's reception and processing capacity, consistent with existing EU allocation proposals and policy debates (\citeproc{ref-thielemann2018refugee}{Thielemann, 2018}). These indicators also provide respondents with clear and easily interpretable signals of a country's size and wealth, which, as research shows, are commonly used when evaluating ``fairness'' in international sharing (\citeproc{ref-kleider2019politics}{Kleider and Stoeckel, 2019}). Even so, respondents may interpret the proportionality by GDP or population in different ways (e.g., a proxy for housing capacity or labour market conditions), which may conflate the meaning of ``capacity'' and limit the interpretation of the results.}

A few observations can be made. First, historical asylum application rates vary substantially across countries, ranging from 16 per 100,000 inhabitants in Poland to 397 per 100,000 in Greece. Second, proportional allocation systems would significantly alter the current distribution of asylum applications. For example, under a population-proportional allocation scheme, Poland would receive an additional 153 applications per 100,000 inhabitants, while Greece would ``send'' 228 applications per 100,000 inhabitants to other countries. We refer to countries in this position as ``net receivers'' or ``net senders.''

Third, proportional allocation rules would generate pronounced cross-country differences depending on the rule being based on population or GDP. For instance, under a GDP-proportional allocation system, Poland would still be a net receiver, but the additional inflow would be considerably smaller, decreasing from 153 applications per 100,000 inhabitants under a population-based system to 76 per 100,000 under a GDP-based system.

\begin{Shaded}
\begin{Highlighting}[]
\NormalTok{tbl\_asylum\_rel }\SpecialCharTok{\%\textgreater{}\%}
\NormalTok{    dplyr}\SpecialCharTok{::}\FunctionTok{mutate}\NormalTok{(}
        \AttributeTok{country =} \FunctionTok{factor}\NormalTok{(}
\NormalTok{            country,}
            \AttributeTok{levels =}\NormalTok{ country[}\FunctionTok{order}\NormalTok{(no\_relocation, }\AttributeTok{decreasing =} \ConstantTok{TRUE}\NormalTok{)]}
\NormalTok{        ),}
        \AttributeTok{d\_pop =}\NormalTok{ population }\SpecialCharTok{{-}}\NormalTok{ no\_relocation,}
        \AttributeTok{d\_gdp =}\NormalTok{ gdp }\SpecialCharTok{{-}}\NormalTok{ no\_relocation}
\NormalTok{    ) }\SpecialCharTok{\%\textgreater{}\%}
\NormalTok{    dplyr}\SpecialCharTok{::}\FunctionTok{arrange}\NormalTok{(}
\NormalTok{        dplyr}\SpecialCharTok{::}\FunctionTok{desc}\NormalTok{(no\_relocation)}
\NormalTok{    ) }\SpecialCharTok{\%\textgreater{}\%}
\NormalTok{    dplyr}\SpecialCharTok{::}\FunctionTok{transmute}\NormalTok{(}
        \AttributeTok{Country =}\NormalTok{ country,}
        \StringTok{"No relocation"} \OtherTok{=}\NormalTok{ no\_relocation,}
        \StringTok{"By population"} \OtherTok{=}\NormalTok{ population,}
        \StringTok{"By GDP"} \OtherTok{=}\NormalTok{ gdp,}
        \StringTok{\textasciigrave{}}\AttributeTok{Δ Population}\StringTok{\textasciigrave{}} \OtherTok{=}\NormalTok{ d\_pop,}
        \StringTok{\textasciigrave{}}\AttributeTok{Δ GDP}\StringTok{\textasciigrave{}} \OtherTok{=}\NormalTok{ d\_gdp}
\NormalTok{    ) }\SpecialCharTok{\%\textgreater{}\%}
\NormalTok{    kableExtra}\SpecialCharTok{::}\FunctionTok{kbl}\NormalTok{(}
        \AttributeTok{caption =} \FunctionTok{paste}\NormalTok{(}
            \StringTok{"Expected asylum applications under alternative"}\NormalTok{,}
            \StringTok{"burden{-}sharing rules (per 100,000 inhabitants)."}
\NormalTok{        ),}
        \AttributeTok{booktabs =} \ConstantTok{TRUE}\NormalTok{,}
        \AttributeTok{linesep =} \StringTok{""}\NormalTok{,}
        \AttributeTok{align =} \FunctionTok{c}\NormalTok{(}\StringTok{"l"}\NormalTok{, }\FunctionTok{rep}\NormalTok{(}\StringTok{"r"}\NormalTok{, }\DecValTok{5}\NormalTok{)),}
        \AttributeTok{escape =} \ConstantTok{FALSE}
\NormalTok{    ) }\SpecialCharTok{\%\textgreater{}\%}
\NormalTok{    kableExtra}\SpecialCharTok{::}\FunctionTok{kable\_styling}\NormalTok{(}
        \AttributeTok{latex\_options =} \FunctionTok{c}\NormalTok{(}\StringTok{"hold\_position"}\NormalTok{),}
        \AttributeTok{full\_width =} \ConstantTok{FALSE}\NormalTok{,}
        \AttributeTok{position =} \StringTok{"center"}\NormalTok{,}
        \AttributeTok{font\_size =} \DecValTok{10}
\NormalTok{    ) }\SpecialCharTok{\%\textgreater{}\%}
\NormalTok{    kableExtra}\SpecialCharTok{::}\FunctionTok{row\_spec}\NormalTok{(}\DecValTok{0}\NormalTok{, }\AttributeTok{bold =} \ConstantTok{TRUE}\NormalTok{) }\SpecialCharTok{\%\textgreater{}\%}
\NormalTok{    kableExtra}\SpecialCharTok{::}\FunctionTok{add\_header\_above}\NormalTok{(}
        \FunctionTok{c}\NormalTok{(}\StringTok{" "} \OtherTok{=} \DecValTok{1}\NormalTok{, }\StringTok{"Levels"} \OtherTok{=} \DecValTok{3}\NormalTok{, }\StringTok{"Differences vs. no relocation"} \OtherTok{=} \DecValTok{2}\NormalTok{),}
        \AttributeTok{bold =} \ConstantTok{TRUE}
\NormalTok{    ) }\SpecialCharTok{\%\textgreater{}\%}
\NormalTok{    kableExtra}\SpecialCharTok{::}\FunctionTok{footnote}\NormalTok{(}
        \AttributeTok{general =} \StringTok{"Entries report expected asylum applications per 100,000 inhabitants."}\NormalTok{,}
        \AttributeTok{threeparttable =} \ConstantTok{TRUE}\NormalTok{,}
        \AttributeTok{general\_title =} \StringTok{"Note:"}
\NormalTok{    )}
\end{Highlighting}
\end{Shaded}

\begin{table}[!h]
\centering
\caption{\label{tab:relocation}Expected asylum applications under alternative burden-sharing rules (per 100,000 inhabitants).}
\centering
\fontsize{10}{12}\selectfont
\begin{threeparttable}
\begin{tabular}[t]{lrrrrr}
\toprule
\multicolumn{1}{c}{\textbf{ }} & \multicolumn{3}{c}{\textbf{Levels}} & \multicolumn{2}{c}{\textbf{Differences vs. no relocation}} \\
\cmidrule(l{3pt}r{3pt}){2-4} \cmidrule(l{3pt}r{3pt}){5-6}
\textbf{Country} & \textbf{No relocation} & \textbf{By population} & \textbf{By GDP} & \textbf{Δ Population} & \textbf{Δ GDP}\\
\midrule
Greece & 397 & 169 & 95 & -228 & -302\\
Sweden & 344 & 169 & 230 & -175 & -114\\
Germany & 313 & 169 & 219 & -144 & -94\\
Bulgaria & 178 & 169 & 65 & -9 & -113\\
France & 151 & 169 & 184 & 18 & 33\\
Spain & 138 & 169 & 136 & 31 & -2\\
Italy & 128 & 169 & 158 & 41 & 30\\
Poland & 16 & 169 & 92 & 153 & 76\\
\bottomrule
\end{tabular}
\begin{tablenotes}
\item \textit{Note:} 
\item Entries report expected asylum applications per 100,000 inhabitants.
\end{tablenotes}
\end{threeparttable}
\end{table}

Finally, the survey included several modules that asked respondents about their socio-demographics, perceptions of the EU asylum system, support for different migration policies, and other background questions. The full questionnaire is in the Appendix. A descriptive analysis of the responses is presented in the technical report by Blasco et al. (\citeproc{ref-blasco2025fairness}{2025}). In our analysis, we included control variables such as trust or political affiliation, as these are associated with attitudes towards migration (e.g., \citeproc{ref-canetti2016threatened}{Canetti et al., 2016}; \citeproc{ref-mitchell2021social}{Mitchell, 2021}). We also consider various socio-demographic factors, such as income and employment status, gender, age, and the parents' country of birth. Whether the respondent was born abroad may be relevant because foreigners may hold different opinions on migration issues than natives (\citeproc{ref-berlinschi2019migrants}{Berlinschi and Harutyunyan, 2019}).

\subsection{Data}\label{data}

Table \ref{tab:assignment-summary} in the Appendix illustrates the random assignment to the three information conditions in each country and overall. Each treatment group has about one third of total respondents per country, reflecting the stratified randomisation design.

Table \ref{tab:summary} illustrates the distribution of key pre-treatment characteristics, including sex, age, education, trust, and perceptions across the three treatment groups. On average, respondents were 50 years old (with 16 years SD), with about half (47\%) identifying as male. For simplicity of exposition and comparability across countries, education was merged into three broad categories: high (ISCED 5-8; 32\% of respondents), medium (ISCED 3-4; 48\%), and low (ISCED 0-2; 19\%). On measures of trust, respondents reported moderate levels of general trust in others (M = 4.5, SD = 2.7), trust in the EU (M = 4.3, SD = 2.7), and trust in government (M = 3.7, SD = 2.9).

Knowledge of the EU asylum system averaged 2.0 (SD = 0.7) on a 1-4 scale, thus the majority expresses no or basic knowledge. Most respondents reported there were ``too many'' asylum seekers in their country, with a mean of 2.6 on 1-3 scale from 1 = too few and 3 = too many. Similarly, on average, respondents view an unbalanced distribution of asylum seekers in the EU, reporting an average score of 3.1 (SD = 2.4) on 0-10 scale for the question about evenness across the EU. Yet, respondents reported broad consensus on the support on the importance of burden sharing among countries, providing equal rights, equal treatment, and having an equal distribution of asylum seekers relatively high, with means ranging from 6.9 to 7.4 on 0-10 scales.

Table \ref{tab:summary} also reports mean differences between the control group and the other treatment groups, with the standardised difference in parenthesis (the difference divided by the pooled standard deviation). For all these variables, the standardised differences between control and other groups is small, suggesting the distributions overlap well across treatment groups.\footnote{
  These differences were estimated using OLS controlling for country fixed effets, thus accounting for the stratified randomisation.}

\begin{Shaded}
\begin{Highlighting}[]
\CommentTok{\# {-}{-}{-}{-}{-}}
\CommentTok{\# 1. Compute SMD}
\CommentTok{\# {-}{-}{-}{-}{-}}
\NormalTok{result }\OtherTok{\textless{}{-}} \FunctionTok{lapply}\NormalTok{(}\FunctionTok{names}\NormalTok{(ds), }\ControlFlowTok{function}\NormalTok{(x) \{}
\NormalTok{    column }\OtherTok{\textless{}{-}}\NormalTok{ ds[[x]]}

    \ControlFlowTok{if}\NormalTok{ (}\SpecialCharTok{!}\FunctionTok{is.numeric}\NormalTok{(column) }\SpecialCharTok{||} \FunctionTok{all}\NormalTok{(}\FunctionTok{is.na}\NormalTok{(column))) \{}
        \FunctionTok{return}\NormalTok{(}\ConstantTok{NULL}\NormalTok{)}
\NormalTok{    \}}

\NormalTok{    formula }\OtherTok{\textless{}{-}} \FunctionTok{reformulate}\NormalTok{(}\FunctionTok{c}\NormalTok{(}\StringTok{"relocation\_treatment"}\NormalTok{, }\StringTok{"country"}\NormalTok{), x)}
\NormalTok{    fit }\OtherTok{\textless{}{-}} \FunctionTok{lm}\NormalTok{(formula, }\AttributeTok{data =}\NormalTok{ ds, }\AttributeTok{contrasts =} \FunctionTok{list}\NormalTok{(}\AttributeTok{country =} \StringTok{"contr.sum"}\NormalTok{))}

    \DocumentationTok{\#\# compute pooled SD}
\NormalTok{    group\_sds }\OtherTok{\textless{}{-}} \FunctionTok{tapply}\NormalTok{(column, ds}\SpecialCharTok{$}\NormalTok{relocation\_treatment, sd, }\AttributeTok{na.rm =} \ConstantTok{TRUE}\NormalTok{)}

\NormalTok{    broom}\SpecialCharTok{::}\FunctionTok{tidy}\NormalTok{(fit) }\SpecialCharTok{|\textgreater{}}
\NormalTok{        dplyr}\SpecialCharTok{::}\FunctionTok{mutate}\NormalTok{(}
            \AttributeTok{pooled\_sd =} \FunctionTok{sqrt}\NormalTok{(}\FunctionTok{sum}\NormalTok{(group\_sds}\SpecialCharTok{\^{}}\DecValTok{2}\NormalTok{) }\SpecialCharTok{/} \FunctionTok{length}\NormalTok{(group\_sds)),}
            \AttributeTok{smd =}\NormalTok{ estimate }\SpecialCharTok{/}\NormalTok{ pooled\_sd,}
            \AttributeTok{variable =}\NormalTok{ x}
\NormalTok{        )}
\NormalTok{\})}

\NormalTok{result }\OtherTok{\textless{}{-}}\NormalTok{ dplyr}\SpecialCharTok{::}\FunctionTok{bind\_rows}\NormalTok{(result)}

\CommentTok{\# {-}{-}{-}{-}{-}}
\CommentTok{\# 2. Build table}
\CommentTok{\# {-}{-}{-}{-}{-}}
\NormalTok{tbl\_covariates }\OtherTok{\textless{}{-}}\NormalTok{ result }\SpecialCharTok{\%\textgreater{}\%}
    \FunctionTok{filter}\NormalTok{(}
        \FunctionTok{str\_detect}\NormalTok{(term, }\StringTok{"relocation\_treatment"}\NormalTok{),}
        \SpecialCharTok{!}\FunctionTok{str\_detect}\NormalTok{(variable, }\StringTok{"weight|restart\_count|\^{}t[DQ]|\^{}tIn"}\NormalTok{),}
\NormalTok{    ) }\SpecialCharTok{\%\textgreater{}\%}
    \FunctionTok{mutate}\NormalTok{(}
        \AttributeTok{stat =} \FunctionTok{sprintf}\NormalTok{(}\StringTok{"\%.2f [\%.2f]"}\NormalTok{, estimate, smd)}
\NormalTok{    ) }\SpecialCharTok{\%\textgreater{}\%}
    \FunctionTok{select}\NormalTok{(variable, term, stat) }\SpecialCharTok{\%\textgreater{}\%}
    \FunctionTok{pivot\_wider}\NormalTok{(}\AttributeTok{names\_from =}\NormalTok{ term, }\AttributeTok{values\_from =}\NormalTok{ stat) }\SpecialCharTok{\%\textgreater{}\%}
    \FunctionTok{mutate}\NormalTok{(}
        \AttributeTok{variable =} \FunctionTok{recode}\NormalTok{(variable,}
            \AttributeTok{age =} \StringTok{"Age (years)"}\NormalTok{,}
            \AttributeTok{male =} \StringTok{"Male (\%)"}\NormalTok{,}
            \AttributeTok{educ\_high =} \StringTok{"High education (\%)"}\NormalTok{,}
            \AttributeTok{educ\_mid =} \StringTok{"Mid education (\%)"}\NormalTok{,}
            \AttributeTok{educ\_low =} \StringTok{"Low education (\%)"}\NormalTok{,}
            \AttributeTok{trust =} \StringTok{"Trust others (0{-}10)"}\NormalTok{,}
            \AttributeTok{trust\_eu =} \StringTok{"Trust in EU (0{-}10)"}\NormalTok{,}
            \AttributeTok{trust\_gov =} \StringTok{"Trust in government (0{-}10)"}\NormalTok{,}
            \AttributeTok{political\_right =} \StringTok{"Political right (0{-}10)"}\NormalTok{,}
            \AttributeTok{share\_legal\_migrants =} \StringTok{"Share legal migrants"}\NormalTok{,}
            \AttributeTok{asylum\_knowledge =} \StringTok{"Knowledge of asylum system (1{-}4)"}\NormalTok{,}
            \AttributeTok{asylum\_fraud\_detected =} \StringTok{"Fraud detection ability (0{-}10)"}\NormalTok{,}
            \AttributeTok{asylum\_decisions\_objective =} \StringTok{"Asylum decisions are objective (0{-}10)"}\NormalTok{,}
            \AttributeTok{asylum\_distr\_even =} \StringTok{"Asylums in EU are even (0{-}10)"}\NormalTok{,}
            \AttributeTok{fair\_share =}  \StringTok{"Asylums in own country (1 = too few, 3 = too many)"}\NormalTok{,}
            \AttributeTok{burden\_sharing\_important =} \StringTok{"Importance of burden sharing (0{-}10)"}\NormalTok{,}
            \AttributeTok{equal\_distribution\_important =} \StringTok{"Importance of equal distribution (0{-}10)"}\NormalTok{,}
            \AttributeTok{equal\_treatment\_important =}  \StringTok{"Importance of equal treatment (0{-}10)"}\NormalTok{,}
            \AttributeTok{equal\_rights\_important =} \StringTok{"Importance of equal rights (0{-}10)"}\NormalTok{,}
            \AttributeTok{quick\_processing\_important =} \StringTok{"Importance of quick process (0{-}10)"}\NormalTok{,}
            \AttributeTok{Duration =} \StringTok{"Completion time (mins)"}\NormalTok{,}
            \AttributeTok{number\_accepted\_refugees =} \StringTok{"Accepted refugees in own country"}\NormalTok{,}
            \AttributeTok{.default =}\NormalTok{ variable}
\NormalTok{        ),}
        \AttributeTok{group =} \FunctionTok{case\_when}\NormalTok{(}
\NormalTok{            variable }\SpecialCharTok{\%in\%} \FunctionTok{c}\NormalTok{(}
                \StringTok{"Age (years)"}\NormalTok{, }\StringTok{"Male (\%)"}\NormalTok{,}
                \StringTok{"High education (\%)"}\NormalTok{, }\StringTok{"Mid education (\%)"}\NormalTok{, }\StringTok{"Low education (\%)"}
\NormalTok{            ) }\SpecialCharTok{\textasciitilde{}} \StringTok{"Sociodemographics"}\NormalTok{,}
\NormalTok{            variable }\SpecialCharTok{\%in\%} \FunctionTok{c}\NormalTok{(}
                \StringTok{"Trust others (0{-}10)"}\NormalTok{, }\StringTok{"Trust in EU (0{-}10)"}\NormalTok{,}
                \StringTok{"Trust in government (0{-}10)"}\NormalTok{, }\StringTok{"Political right (0{-}10)"}
\NormalTok{            ) }\SpecialCharTok{\textasciitilde{}} \StringTok{"Political attitudes"}\NormalTok{,}
\NormalTok{            variable }\SpecialCharTok{\%in\%} \FunctionTok{c}\NormalTok{(}
                \StringTok{"Knowledge of asylum system (1{-}4)"}\NormalTok{,}
                \StringTok{"Fraud detection ability (0{-}10)"}
\NormalTok{            ) }\SpecialCharTok{\textasciitilde{}} \StringTok{"Knowledge"}\NormalTok{,}
\NormalTok{            variable }\SpecialCharTok{\%in\%} \StringTok{"Completion time (mins)"} \SpecialCharTok{\textasciitilde{}} \StringTok{"Survey administration"}\NormalTok{,}
            \ConstantTok{TRUE} \SpecialCharTok{\textasciitilde{}} \StringTok{"Other"}
\NormalTok{        )}
\NormalTok{    ) }\SpecialCharTok{\%\textgreater{}\%}
    \FunctionTok{arrange}\NormalTok{(}\FunctionTok{desc}\NormalTok{(group)) }\SpecialCharTok{\%\textgreater{}\%} 
    \FunctionTok{mutate}\NormalTok{(}
        \AttributeTok{variable =} \FunctionTok{str\_replace}\NormalTok{(variable, }\StringTok{"Importance of"}\NormalTok{, }\StringTok{"Value of"}\NormalTok{),}
\NormalTok{    ) }

\NormalTok{groups }\OtherTok{\textless{}{-}} \FunctionTok{unique}\NormalTok{(tbl\_covariates}\SpecialCharTok{$}\NormalTok{group)}

\NormalTok{tbl\_covariates }\SpecialCharTok{\%\textgreater{}\%}
    \FunctionTok{mutate}\NormalTok{(}\AttributeTok{variable =} \FunctionTok{gsub}\NormalTok{(}\StringTok{"\_"}\NormalTok{, }\StringTok{" "}\NormalTok{, variable)) }\SpecialCharTok{\%\textgreater{}\%}
    \FunctionTok{select}\NormalTok{(}
        \AttributeTok{Variable =}\NormalTok{ variable,}
        \StringTok{\textasciigrave{}}\AttributeTok{Relative treatment}\StringTok{\textasciigrave{}} \OtherTok{=} \FunctionTok{contains}\NormalTok{(}\StringTok{"Relative"}\NormalTok{),}
        \StringTok{\textasciigrave{}}\AttributeTok{Absolute treatment}\StringTok{\textasciigrave{}} \OtherTok{=} \FunctionTok{contains}\NormalTok{(}\StringTok{"Absolute"}\NormalTok{)}
\NormalTok{    ) }\SpecialCharTok{\%\textgreater{}\%} 
\NormalTok{    kableExtra}\SpecialCharTok{::}\FunctionTok{kbl}\NormalTok{(}
        \AttributeTok{booktabs =} \ConstantTok{TRUE}\NormalTok{,}
        \AttributeTok{align =} \FunctionTok{c}\NormalTok{(}\StringTok{"l"}\NormalTok{, }\StringTok{"c"}\NormalTok{, }\StringTok{"c"}\NormalTok{),}
        \AttributeTok{escape =} \ConstantTok{TRUE}\NormalTok{,}
        \AttributeTok{caption =} \FunctionTok{paste}\NormalTok{(}
            \StringTok{"Covariate balance across treatment arms."}\NormalTok{,}
            \StringTok{"Mean differences with respect to the control group."}
\NormalTok{        )}
\NormalTok{    ) }\SpecialCharTok{\%\textgreater{}\%}
\NormalTok{    kableExtra}\SpecialCharTok{::}\FunctionTok{kable\_styling}\NormalTok{(}
        \AttributeTok{full\_width =} \ConstantTok{FALSE}\NormalTok{, }
        \AttributeTok{position =} \StringTok{"center"}
\NormalTok{    ) }\SpecialCharTok{\%\textgreater{}\%}
\NormalTok{    kableExtra}\SpecialCharTok{::}\FunctionTok{pack\_rows}\NormalTok{(groups[}\DecValTok{1}\NormalTok{],}
        \AttributeTok{start\_row =} \DecValTok{1}\NormalTok{,}
        \AttributeTok{end\_row =} \FunctionTok{sum}\NormalTok{(tbl\_covariates}\SpecialCharTok{$}\NormalTok{group }\SpecialCharTok{\%in\%}\NormalTok{ groups[}\DecValTok{1}\NormalTok{])}
\NormalTok{    ) }\SpecialCharTok{\%\textgreater{}\%}
\NormalTok{    kableExtra}\SpecialCharTok{::}\FunctionTok{pack\_rows}\NormalTok{(groups[}\DecValTok{2}\NormalTok{],}
        \AttributeTok{start\_row =} \FunctionTok{sum}\NormalTok{(tbl\_covariates}\SpecialCharTok{$}\NormalTok{group }\SpecialCharTok{\%in\%}\NormalTok{ groups[}\DecValTok{1}\NormalTok{]) }\SpecialCharTok{+} \DecValTok{1}\NormalTok{,}
        \AttributeTok{end\_row =} \FunctionTok{sum}\NormalTok{(tbl\_covariates}\SpecialCharTok{$}\NormalTok{group }\SpecialCharTok{\%in\%}\NormalTok{ groups[}\DecValTok{1}\SpecialCharTok{:}\DecValTok{2}\NormalTok{])}
\NormalTok{    ) }\SpecialCharTok{\%\textgreater{}\%}
\NormalTok{    kableExtra}\SpecialCharTok{::}\FunctionTok{pack\_rows}\NormalTok{(groups[}\DecValTok{3}\NormalTok{],}
        \AttributeTok{start\_row =} \FunctionTok{sum}\NormalTok{(tbl\_covariates}\SpecialCharTok{$}\NormalTok{group }\SpecialCharTok{\%in\%}\NormalTok{ groups[}\DecValTok{1}\SpecialCharTok{:}\DecValTok{2}\NormalTok{]) }\SpecialCharTok{+} \DecValTok{1}\NormalTok{,}
        \AttributeTok{end\_row =} \FunctionTok{sum}\NormalTok{(tbl\_covariates}\SpecialCharTok{$}\NormalTok{group }\SpecialCharTok{\%in\%}\NormalTok{ groups[}\DecValTok{1}\SpecialCharTok{:}\DecValTok{3}\NormalTok{])}
\NormalTok{    ) }\SpecialCharTok{\%\textgreater{}\%} 
\NormalTok{    kableExtra}\SpecialCharTok{::}\FunctionTok{pack\_rows}\NormalTok{(groups[}\DecValTok{4}\NormalTok{],}
        \AttributeTok{start\_row =} \FunctionTok{sum}\NormalTok{(tbl\_covariates}\SpecialCharTok{$}\NormalTok{group }\SpecialCharTok{\%in\%}\NormalTok{ groups[}\DecValTok{1}\SpecialCharTok{:}\DecValTok{3}\NormalTok{]) }\SpecialCharTok{+} \DecValTok{1}\NormalTok{,}
        \AttributeTok{end\_row =} \FunctionTok{nrow}\NormalTok{(tbl\_covariates)}
\NormalTok{    ) }\SpecialCharTok{\%\textgreater{}\%} 
\NormalTok{    kableExtra}\SpecialCharTok{::}\FunctionTok{footnote}\NormalTok{(}
        \AttributeTok{general =} \FunctionTok{paste}\NormalTok{(}
            \StringTok{"Entries report OLS coefficients from regressions"}\NormalTok{, }
            \StringTok{"of each covariate on treatment indicators and country fixed effects."}\NormalTok{,}
            \StringTok{"Standardized mean differences (SMDs) are reported in brackets."}
\NormalTok{        ),}
        \AttributeTok{threeparttable =} \ConstantTok{TRUE}\NormalTok{,}
        \AttributeTok{general\_title =} \StringTok{"Note:"}
\NormalTok{    )}
\end{Highlighting}
\end{Shaded}

\begin{table}
\centering
\caption{\label{tab:summary}Covariate balance across treatment arms. Mean differences with respect to the control group.}
\centering
\begin{threeparttable}
\begin{tabular}[t]{lcc}
\toprule
Variable & Relative treatment & Absolute treatment\\
\midrule
\addlinespace[0.3em]
\multicolumn{3}{l}{\textbf{Survey administration}}\\
\hspace{1em}Completion time (mins) & -1.11 [-0.06] & -0.25 [-0.01]\\
\addlinespace[0.3em]
\multicolumn{3}{l}{\textbf{Sociodemographics}}\\
\hspace{1em}Age (years) & -0.28 [-0.02] & -0.15 [-0.01]\\
\hspace{1em}Male (\%) & 0.01 [0.03] & 0.02 [0.03]\\
\hspace{1em}High education (\%) & -0.01 [-0.02] & 0.01 [0.02]\\
\hspace{1em}Mid education (\%) & 0.00 [0.00] & -0.01 [-0.01]\\
\hspace{1em}Low education (\%) & 0.01 [0.03] & -0.00 [-0.00]\\
\addlinespace[0.3em]
\multicolumn{3}{l}{\textbf{Political attitudes}}\\
\hspace{1em}Trust others (0-10) & 0.07 [0.03] & 0.09 [0.03]\\
\hspace{1em}Trust in government (0-10) & 0.16 [0.06] & 0.17 [0.06]\\
\hspace{1em}Trust in EU (0-10) & 0.12 [0.04] & 0.12 [0.04]\\
\hspace{1em}Political right (0-10) & -0.07 [-0.03] & -0.07 [-0.03]\\
\addlinespace[0.3em]
\multicolumn{3}{l}{\textbf{Other}}\\
\hspace{1em}Asylums in EU are even (0-10) & -0.10 [-0.04] & -0.01 [-0.01]\\
\hspace{1em}Value of burden sharing (0-10) & -0.03 [-0.01] & 0.03 [0.01]\\
\hspace{1em}Value of equal distribution (0-10) & -0.04 [-0.02] & -0.02 [-0.01]\\
\hspace{1em}Value of equal rights (0-10) & -0.01 [-0.00] & 0.07 [0.03]\\
\hspace{1em}Value of equal treatment (0-10) & -0.07 [-0.03] & 0.04 [0.01]\\
\hspace{1em}Value of quick process (0-10) & -0.09 [-0.04] & 0.02 [0.01]\\
\hspace{1em}Accepted refugees in own country & -65.91 [-0.00] & 1159.07 [0.04]\\
\hspace{1em}Asylum decisions are objective (0-10) & -0.06 [-0.02] & -0.11 [-0.04]\\
\hspace{1em}Share legal migrants & -0.20 [-0.01] & 0.81 [0.03]\\
\hspace{1em}Fraud detection ability (0-10) & 0.09 [0.03] & 0.12 [0.04]\\
\bottomrule
\end{tabular}
\begin{tablenotes}
\item \textit{Note:} 
\item Entries report OLS coefficients from regressions of each covariate on treatment indicators and country fixed effects. Standardized mean differences (SMDs) are reported in brackets.
\end{tablenotes}
\end{threeparttable}
\end{table}

\section{Discussion}\label{discussion}

This study used an experimental survey to examine Europeans' perceptions of the fairness of different mechanisms for distributing asylum applications among EU countries. By varying the information provided about the consequences of these mechanisms, it explored the extent to which public perceptions of fairness are driven by concerns about outcomes versus principles of fairness.

Consistent with prior research (\citeproc{ref-bansak2017europeans}{Bansak et al., 2017}), we find that a substantial majority (76\%) of respondents across eight EU countries view some form of proportional relocation of applicants as \emph{fairer} than no relocation at all. Moreover, most people view proportional relocation systems as \emph{fairer} even when it would result in their own country receiving more asylum applications, as the share of those ranking relocation as fairer than no relocation exceeds 50\% in all countries.

Yet, we find substantial variation in fairness perceptions across countries. In the control group, where no information about the outcomes was provided, we find that between 58\% respondents in Germany to 87\% respondents in Greece view some form of within-EU relocation as fairer than no relocation at all. This gap could be partly explained by concerns about the outcomes. For instance, a border country like Greece would experience more applications under a no relocation system, while Germany would experience fewer applications. This finding underscores the importance of disentangling concerns about outcomes from fairness.

Bansak et al. (\citeproc{ref-bansak2017europeans}{2017}) finds that varying the information available about outcomes significantly affects public support for different relocation systems. Building on this insight, our research demonstrates that providing information about the consequences is key in shaping public perceptions of the fairness of these systems. Thus, our study establishes a direct link between political support identified by Bansak et al. (\citeproc{ref-bansak2017europeans}{2017}) and the perceived fairness of different policy options. As such, our study adds to a growing literature exploring the role of public perceptions of fairness in regulation, such as Brick and Visser (\citeproc{ref-brick2015fair}{2015}) for climate regulation.

By providing information about the consequences of the allocation schemes, we find that when individuals learn about the outcomes, the initial preference for proportional schemes is offset if the relocation mechanism results in one standard deviation more asylum applications. In other words, Europeans are more likely to perceive proportional allocation systems as fair when the increase in the number of asylum seekers implied by these systems is not too large. A first key takeaway is that information matters: receiving information about the consequences of proportional mechanisms may alter perceptions of fairness, ultimately shaping public support for relocation policies.

Another key implication is that framing matters. Differences that appear large in absolute terms may seem small when presented in relative terms. Indeed, in our randomised experiment, we find that expressing the outcomes of different relocation systems in absolute numbers leads some individuals to react more strongly to the information. When no information is provided, a regime that increases expected asylum applications by one standard deviation is 13\% less likely to be seen as fair. Once information is revealed, fairness drops by 40\% with relative information and by 56\% with absolute information. Thus, providing absolute information produces an additional 16-percentage point decrease in perceived fairness compared to relative information. This finding highlights another way in which public opinion can be influenced by the framing of identical information.

Our analysis of heterogeneous effects provides additional evidence on the mechanisms through which absolute vs relative information works. These differing impacts are not uniform across the population but are highly dependent on individual factors, including political orientation and demographics. Across most education groups, we find a consistently negative and statistically significant correlation between the number of asylum seekers and perceived fairness. The magnitude of this association was larger in the absolute-information group, moderate in the relative-information group, and modest in the control group. However, the effect of asylum seekers was weakest for low-educated respondents in all groups. These findings suggest a minor sensitivity of low-educated individuals to information. With regard to political orientation, both GDP- and population-based relocation schemes were consistently viewed as fairer than no relocation, with the largest effects observed for left-leaning respondents under absolute and control conditions.

Our research relates to conspicuous experimental literature in other fields that has shown the sensitivity of human decision-making to absolute vs relative values (\citeproc{ref-kratochvil2024fallacy}{Kratochvı́l and Havlı́ček, 2024}; \citeproc{ref-malenka1993framing}{Malenka et al., 1993}). In medical research, for example, effect sizes or risk reductions of interventions may appear very large or quite small, depending on the use of relative or absolute measures on the same dataset (\citeproc{ref-malenka1993framing}{Malenka et al., 1993}). The use of relative risk reduction has been shown to be more likely to influence the decision to accept a treatment than absolute risk reduction. In political communication, the main reason why relative values are so widespread is that absolute values are mostly not of much use for comparison among countries because they vary in population size. It is worth noticing, however, that the use of relative information may lead to misleading conclusions (\citeproc{ref-silva2020per}{Silva, 2020}). They should, therefore, be used when the measure of interest increases proportionally with population size. Our study provides additional evidence on the effect of relative vs absolute communication in the field of migration policies.

This study has some key limitations. As the evidence is survey-based, the findings rely on self-reported attitudes and are susceptible to biases such as social desirability or a gap between stated and actual behaviour. The scope of the research is also geographically constrained: drawing data from eight countries limits the generalizability of the results to global or even broader regional contexts. Finally, the use of hypothetical scenarios, though necessary for experimental control, means that the observed reactions to absolute versus relative information may not perfectly translate to real-world policy decisions.

\section*{References}\label{references}
\addcontentsline{toc}{section}{References}

\phantomsection\label{refs}
\begin{CSLReferences}{1}{0}
\bibitem[\citeproctext]{ref-adams1965inequity}
Adams, J.S., 1965. Inequity in social exchange, in: Advances in Experimental Social Psychology. Elsevier, pp. 267--299.

\bibitem[\citeproctext]{ref-alesina2023immigration}
Alesina, A., Miano, A., Stantcheva, S., 2023. Immigration and redistribution. The Review of Economic Studies 90, 1--39.

\bibitem[\citeproctext]{ref-almaas2010fairness}
Almås, I., Cappelen, A.W., Sørensen, E.Ø., Tungodden, B., 2010. Fairness and the development of inequality acceptance. Science 328, 1176--1178.

\bibitem[\citeproctext]{ref-babcock1995biased}
Babcock, L., Loewenstein, G., Issacharoff, S., Camerer, C., 1995. Biased judgments of fairness in bargaining. The American Economic Review 85, 1337--1343.

\bibitem[\citeproctext]{ref-bansak2017europeans}
Bansak, K., Hainmueller, J., Hangartner, D., 2017. Europeans support a proportional allocation of asylum seekers. Nature Human Behaviour 1, 0133.

\bibitem[\citeproctext]{ref-berlinschi2019migrants}
Berlinschi, R., Harutyunyan, A., 2019. Do migrants think differently? Evidence from eastern european and post-soviet states. International Migration Review 53, 831--868.

\bibitem[\citeproctext]{ref-blasco2025fairness}
Blasco, A., Cabeza Martinez, B., Icardi, R., Krawczyk, M., Seiger, F., 2025. Public perceptions of fairness in the european migration and asylum system. Publications Office of the European Union, Luxembourg. \url{https://doi.org/10.2760/1271462}

\bibitem[\citeproctext]{ref-brick2015fair}
Brick, K., Visser, M., 2015. What is fair? An experimental guide to climate negotiations. European Economic Review 74, 79--95.

\bibitem[\citeproctext]{ref-canetti2016threatened}
Canetti, D., Snider, K.L., Pedersen, A., Hall, B.J., 2016. Threatened or threatening? How ideology shapes asylum seekers' immigration policy attitudes in israel and australia. Journal of refugee studies 29, 583--606.

\bibitem[\citeproctext]{ref-cappelen2023second}
Cappelen, A.W., Cappelen, C., Tungodden, B., 2023. Second-best fairness: The trade-off between false positives and false negatives. American Economic Review 113, 2458--2485.

\bibitem[\citeproctext]{ref-cappelen2023trade}
Cappelen, C.W., Sicakkan, H.G., Van Wolleghem, P.G., 2023. The trade-off between admitting and paying: Experimental evidence on attitudes towards asylum responsibility-sharing. European Union Politics 24, 470--493.

\bibitem[\citeproctext]{ref-carattini2019win}
Carattini, S., Kallbekken, S., Orlov, A., 2019. How to win public support for a global carbon tax. Nature 565, 289--291.

\bibitem[\citeproctext]{ref-corneo2002individual}
Corneo, G., Grüner, H.P., 2002. Individual preferences for political redistribution. Journal of public Economics 83, 83--107.

\bibitem[\citeproctext]{ref-corre2024extent}
Corre, T.L., Tilley, J., 2024. To what extent does asylum policy match public policy preferences? International Migration Review 01979183241253502.

\bibitem[\citeproctext]{ref-dennison2020basic}
Dennison, J., 2020. A basic human values approach to migration policy communication. Data \& Policy 2, e18.

\bibitem[\citeproctext]{ref-dinas2019waking}
Dinas, E., Matakos, K., Xefteris, D., Hangartner, D., 2019. Waking up the golden dawn: Does exposure to the refugee crisis increase support for extreme-right parties? Political analysis 27, 244--254.

\bibitem[\citeproctext]{ref-durante2014preferences}
Durante, R., Putterman, L., Van der Weele, J., 2014. Preferences for redistribution and perception of fairness: An experimental study. Journal of the European Economic Association 12, 1059--1086.

\bibitem[\citeproctext]{ref-eu2024pact}
European Commission, 2024. \href{https://op.europa.eu/en/publication-detail/-/publication/a426d987-0767-11ef-a251-01aa75ed71a1/}{{Pact on migration and asylum}}.

\bibitem[\citeproctext]{ref-facchini2022countering}
Facchini, G., Margalit, Y., Nakata, H., 2022. Countering public opposition to immigration: The impact of information campaigns. European Economic Review 141, 103959.

\bibitem[\citeproctext]{ref-fong2001social}
Fong, C., 2001. Social preferences, self-interest, and the demand for redistribution. Journal of Public economics 82, 225--246.

\bibitem[\citeproctext]{ref-fong2011fairness}
Fong, C.M., Luttmer, E.F., 2011. Do fairness and race matter in generosity? Evidence from a nationally representative charity experiment. Journal of Public Economics 95, 372--394.

\bibitem[\citeproctext]{ref-haaland2020labor}
Haaland, I., Roth, C., 2020. Labor market concerns and support for immigration. Journal of Public Economics 191, 104256.

\bibitem[\citeproctext]{ref-hainmueller2015hidden}
Hainmueller, J., Hopkins, D.J., 2015. The hidden american immigration consensus: A conjoint analysis of attitudes toward immigrants. American journal of political science 59, 529--548.

\bibitem[\citeproctext]{ref-halla2017immigration}
Halla, M., Wagner, A.F., Zweimüller, J., 2017. Immigration and voting for the far right. Journal of the European economic association 15, 1341--1385.

\bibitem[\citeproctext]{ref-heizmann2020willing}
Heizmann, B., Ziller, C., 2020. Who is willing to share the burden? Attitudes towards the allocation of asylum seekers in comparative perspective. Social Forces 98, 1026--1051.

\bibitem[\citeproctext]{ref-jeannet2021asylum}
Jeannet, A.-M., Heidland, T., Ruhs, M., 2021. What asylum and refugee policies do europeans want? Evidence from a cross-national conjoint experiment. European Union Politics 22, 353--376.

\bibitem[\citeproctext]{ref-kleider2019politics}
Kleider, H., Stoeckel, F., 2019. The politics of international redistribution: Explaining public support for fiscal transfers in the EU. European Journal of Political Research 58, 4--29.

\bibitem[\citeproctext]{ref-konow2003fairest}
Konow, J., 2003. Which is the fairest one of all? A positive analysis of justice theories. Journal of economic literature 41, 1188--1239.

\bibitem[\citeproctext]{ref-konow2000fair}
Konow, J., 2000. Fair shares: Accountability and cognitive dissonance in allocation decisions. American economic review 90, 1072--1092.

\bibitem[\citeproctext]{ref-kratochvil2024fallacy}
Kratochvı́l, L., Havlı́ček, J., 2024. The fallacy of global comparisons based on per capita measures. Royal Society Open Science 11, 230832.

\bibitem[\citeproctext]{ref-kustov2025immigration}
Kustov, A., Landgrave, M., 2025. Immigration is difficult?! Informing voters about immigration policy fosters pro-immigration views. Journal of Experimental Political Science 1--13.

\bibitem[\citeproctext]{ref-kuziemko2015elastic}
Kuziemko, I., Norton, M.I., Saez, E., Stantcheva, S., 2015. How elastic are preferences for redistribution? Evidence from randomized survey experiments. American Economic Review 105, 1478--1508.

\bibitem[\citeproctext]{ref-maestre2019perceived}
Maestre-Andrés, S., Drews, S., Van Den Bergh, J., 2019. Perceived fairness and public acceptability of carbon pricing: A review of the literature. Climate policy 19, 1186--1204.

\bibitem[\citeproctext]{ref-malenka1993framing}
Malenka, D.J., Baron, J.A., Johansen, S., Wahrenberger, J.W., Ross, J.M., 1993. The framing effect of relative and absolute risk. Journal of general internal medicine 8, 543--548.

\bibitem[\citeproctext]{ref-mitchell2021social}
Mitchell, J., 2021. Social trust and anti-immigrant attitudes in europe: A longitudinal multi-level analysis. Frontiers in sociology 6, 604884.

\bibitem[\citeproctext]{ref-pedersen2017ratio}
Pedersen, R.T., 2017. Ratio bias and policy preferences: How equivalency framing of numbers can affect attitudes. Political Psychology 38, 1103--1120.

\bibitem[\citeproctext]{ref-reyna2008numeracy}
Reyna, V.F., Brainerd, C.J., 2008. Numeracy, ratio bias, and denominator neglect in judgments of risk and probability. Learning and individual differences 18, 89--107.

\bibitem[\citeproctext]{ref-schwartz1992universals}
Schwartz, S.H., 1992. Universals in the content and structure of values: Theoretical advances and empirical tests in 20 countries, in: Advances in Experimental Social Psychology. Elsevier, pp. 1--65.

\bibitem[\citeproctext]{ref-silva2020per}
Silva, W., 2020. Per capita death and infection rates should be avoided in international comparisons. Public Health 186, 18.

\bibitem[\citeproctext]{ref-solodoch2024overburdened}
Solodoch, O., 2024. Overburdened? How refugee dispersal policies can mitigate NIMBYism and public backlash. Journal of European Public Policy 31, 2455--2482.

\bibitem[\citeproctext]{ref-thielemann2018refugee}
Thielemann, E., 2018. Why refugee burden-sharing initiatives fail: Public goods, free-riding and symbolic solidarity in the EU. JCMS: Journal of Common Market Studies 56, 63--82.

\bibitem[\citeproctext]{ref-thielemann2003between}
Thielemann, E.R., 2003. Between interests and norms: Explaining burden-sharing in the european union. Journal of refugee studies 16, 253--273.

\bibitem[\citeproctext]{ref-van2006making}
Van Oorschot, W., 2006. Making the difference in social europe: Deservingness perceptions among citizens of european welfare states. Journal of European social policy 16, 23--42.

\bibitem[\citeproctext]{ref-van2000should}
Van Oorschot, W., 2000. Who should get what, and why? On deservingness criteria and the conditionality of solidarity among the public. Policy \& Politics 28, 33--48.

\bibitem[\citeproctext]{ref-walster1978equity}
Walster, E.H., Walster, G.W., Berscheid, E., Austin, W., Traupmann, J., Utne, M.K., 1978. Equity: Theory and research. Allyn; Bacon.

\end{CSLReferences}

\end{document}
